
\documentclass[12pt]{article}

%%% These are some packages that are useful
\usepackage{amsfonts, lipsum}
\usepackage[scr]{rsfso}
\usepackage{amsmath,amssymb, amscd,amsbsy, bbm, amsthm, enumerate}
\usepackage[top=1in, bottom=1in, left=.45in, right=.45in]{geometry}
\usepackage[unicode]{hyperref}
\usepackage{tikz, xcolor, fancyhdr}
\usepackage{graphicx, caption, subcaption}

%%% Page formatting
%\setlength{\headsep}{30pt}
\setlength{\parindent}{0pt}
\setlength{\textheight}{9in}



%%% These define our question environment and help number things correctly
\theoremstyle{definition}
\newtheorem{thm}{Theorem}
\newtheorem{question}[thm]{Question}
\newtheorem{questionS}[thm]{***Question}
\newtheorem{prop}[thm]{Proposition}
\newtheorem{lem}[thm]{Lemma}
\newtheorem{cor}[thm]{Corollary}
\newtheorem{defn}[thm]{Definition}
\newtheorem{rem}[thm]{Remark}
\newtheorem{ex}[thm]{Example}

%%% These are some shortcuts that are handy
\def\real{{\mathbb R}}
\def\Natural{\mathbb{N}}
\def\dx{\textnormal{dx}}
\def\dy{\textnormal{dy}}
\def\dz{\textnormal{dz}}
\def\dt{\textnormal{dt}}
\def\Re{\textnormal{Re}}
\def\Im{\textnormal{Im}}
\def\exp{\textnormal{exp}}
\def\interior{\textnormal{interior}}
\def\al{\alpha}
\def\del{\delta}
\def\Del{\Delta}
\def\gam{\gamma}
\def\Gam{\Gamma}
\def\Om{\Omega}
\def\ep{\varepsilon}
\def\lam{\lambda}
\def\rational{{\mathbb Q}}
\def\integer{{\mathbb Z}}
\def\Q{{\mathbb Q}}
\def\Z{{\mathbb Z}}
\def\N{{\mathbb N}}
\def\R{{\mathbb R}}
\def\grad{\nabla}
\def\C{\mathcal C}
\def\P{\mathbb{P}}
\def\T{\mathcal T}
\def\I{\mathcal I}
\newcommand{\abs}[1]{\left| #1 \right|}
\newcommand{\inner}[1]{\langle #1 \rangle}
\newcommand{\norm}[1]{\left\lVert#1\right\rVert}
\newcommand{\spanvect}{\textnormal{span}}
\newcommand{\union}{\cup}
\newcommand{\Union}{\bigcup}
\def\intersect{\cap}
\def\Intersect{\bigcap}
\renewcommand{\neg}{\mathord{\sim}}
\newcommand{\ds}{\displaystyle}





%%%%%%%%%%%%%%%%%%%%%%%%%%%%%%%%%
%%%%%%%%%%%%%%%%%%%%%%%%%%%%%%%%%
\newcommand{\name}{Duff Baker-Jarvis}
%%%%%%%%%%%%%%%%%%%%%%%%%%%%%%%%%
%%%%%%%%%%%%%%%%%%%%%%%%%%%%%%%%%
%%% Document Starts now
\begin{document}
\title{Multidimensional Real Analysis - Duistermaat, Kolk}
\date{\today}
\author{Dr. Duff Baker-Jarvis}
\maketitle

\begin{enumerate}
    \item Chapter 1 - Continuity
    \begin{enumerate}
	\item section 1.1 - Inner product and Norm
	    \begin{itemize}
	    \item (Defn) Inner product and Norm
	    \item (Thm) Properties of Inner Product
	    \item (Thm) Cauchy-Schwarz Inequality\\
		Proof Idea: Decompose a vector $ x $ as its projection onto $ y $ plus its orthogonal component. Take the norm squared and solve for the orthogonal component. Its nonegegativity forces the CS inequality.
	    \item (Thm) Inequalities related to the norm
		\begin{itemize}
		    \item Triangle inequality and extended triangle inequality
		    \item Reverse triangle inequality
		    \item $ \ds |x_j|\leq ||x|| \leq \sum_{i=1}^{n} |x_i| \leq \sqrt{n}||x||  $\\
			Proof Idea: First ineqality is straight forward, second inequality can be seen by decomposing the vector and applying the triangle inequality. Third inequality can be seend by forming the vector $ x_{\text{abs}}| $ composed of the absolute values. Then take the inner product with the vector of all 1s and apply C.S.
		    \item $ \max_{j}|x_j| \leq ||x|| \leq \sqrt{n} \max_{j} |x_j|$. This implies that for a cube or side length 2 is enclosed in a sphere of radius $ \sqrt{n} $ which is enclosed in a cube of side length $ 2\sqrt{n} $.\\
			Proof Idea: First inequality follows from the previous part. 2nd inequality can be seen by replacing each component of $ x $ with the maximum component. The norm of the new vector is the RHS.
		\end{itemize}
	    \item (Defn) Distance between two vectors
	    \item (Defn) Convergence or limit of a sequence of vectors in $ \R^n $.
	    \item (Thm) A sequence of vectors converges iff each sequence of components converges.
	    \item (Thm) The limit operation on vectors ist still linear.
	    \end{itemize}
	\item Section 1.2 - Open and closed sets
	    \begin{itemize}
		\item (Defn) Open Ball/neighborhood
	    \end{itemize}
	    
    \end{enumerate}
    
\end{enumerate}


\end{document}
